% !TeX root = ../../Book4.tex
% 纲要标题：### 26.4 非线性修正与边界图册
% 纲要位置：几何学计划纲要.md，第 4320 行。
% 正文状态：未撰写；本文件目前是框架。

% 纲要内容开始
% 
% * 二次余项的 Lipschitz 估计
% * 压缩映射或定量隐函数定理
% * 粘合解的唯一性
% * 所有邻近轨线的反向识别
% * 一维紧化的边界结构
% * 颈长坐标的光滑性
% 
%   证明粘合参数给出真正的端部图册，而不只是一组存在解。
% 
% * 边界覆盖的完备性
% 
%   利用反粘合排除额外端点；这是把计数等式转为 d²=0 的最后分析前提。
% 
% 纲要内容结束

\section{非线性修正与边界图册}
\label{g4:sec:2604}

待撰写

\subsection{二次余项的 Lipschitz 估计}
\label{g4:subsec:2604:01}
待撰写

\subsection{压缩映射或定量隐函数定理}
\label{g4:subsec:2604:02}
待撰写

\subsection{粘合解的唯一性}
\label{g4:subsec:2604:03}
待撰写

\subsection{所有邻近轨线的反向识别}
\label{g4:subsec:2604:04}
待撰写

\subsection{一维紧化的边界结构}
\label{g4:subsec:2604:05}
待撰写

\subsection{颈长坐标的光滑性}
\label{g4:subsec:2604:06}
待撰写

\subsection{边界覆盖的完备性}
\label{g4:subsec:2604:07}
待撰写
% 纲要条目排为编号小节标题，并进入目录与书签；说明段仍保留于纲要注释中。
